% !TEX root = ../b-rg.tex
\chapter{Numerals}\label{ch:numerals}
\stub

% --- Planned sections ---
% \section{Cardinal numerals}
%   zefr, un, dou, tre, catr, cinq, segs, set, ogt, nouf, deg; onç ... novoç;
%   veint, trent, carant ... nonant; cent, mil; murðon 'ten thousand' and
%   cent murðon 'million' -- is the large-number system myriad-based?
% \section{Ordinal numerals}
%   primer, dovem (cf. the 'primer ... dovem' list in Scaunç) ...
% \section{Collective and other derived numerals}
%   doçan 'dozen'; ogtoç 'eighteensome'; carant, cinquant, quinç as
%   'period of N years'; un des y mil 'one in a million'; fractions.
% \section{Numerals in the noun phrase}
%   Position; numerals as number marking; verb agreement.
% \section{Measures, money and dates}
%   Dime units, the gemma, writing dates and clock times (notes 132.5,
%   170.5, 183.5); month and day names are in Lexical Fields.
